\chapter{2006年上海卷（秋理）}
\section{填空题}
\begin{problem}
已知集合 \(A = \{-1, 3, 2m - 1\}\)，集合 \(B = \{3, m^2\}\)，若 \(B \subseteq A\)，则实数 \(m = \)\fillinblank{}.
\end{problem}

\begin{answer}
\(1\)
\end{answer}

\begin{problem}
已知圆 \(x^{2} - 4x - 4 + y^{2} = 0\) 的圆心是点 \(P\)，则点 \(P\) 到直线 \(x - y - 1 = 0\) 的距离是\fillinblank{}.
\end{problem}

\begin{answer}
\(\frac{\sqrt{2}}{2}\)
\end{answer}

\begin{problem}
若函数 \( f(x) = a^x \)（\(a > 0\)，且 \(a \neq 1\)）的反函数的图象过点 \((2,-1)\)，则 \( a = \)\fillinblank{}.
\end{problem}

\begin{answer}
\(\frac{1}{2}\)
\end{answer}

\begin{problem}
计算 \(\lim\limits_{n\to \infty}\frac{\mathrm C_n^3}{n^3 + 1} =\)\fillinblank{}.
\end{problem}

\begin{answer}
\(\frac{1}{6}\)
\end{answer}

\begin{problem}
若复数 \(z\) 同时满足 \(z - \overline{z} = 2\i\)，\(\overline{z} = \i z\) （\(\i\) 为虚数单位），则 \(z =\fillinblank{}\).
\end{problem}

\begin{answer}
\(-1+\i\)
\end{answer}

\begin{problem}
如果 \(\cos \alpha = \frac{1}{5}\)，且 \(\alpha\) 是第四象限的角，那么 \(\cos \left(\alpha + \frac{\pi}{2}\right) =\)\fillinblank{}.
\end{problem}

\begin{answer}
\(\frac{2\sqrt{6}}{5}\)
\end{answer}

\begin{problem}
已知椭圆中心在原点，一个焦点为 \(F(-2\sqrt{3}, 0)\)，且长轴长是短轴长的 2 倍，则该椭圆的标准方程是\fillinblank{}.
\end{problem}

\begin{answer}
\(\frac{x^2}{16}+\frac{y^2}{4}=1\)
\end{answer}

\begin{problem}
在极坐标系中，\(O\) 是极点，设点 \(A\left(4, \frac{\pi}{3}\right)\)，\(B\left(5, -\frac{5\pi}{6}\right)\)，则 \(\triangle OAB\) 的面积是\fillinblank{}.
\end{problem}

\begin{answer}
\(5\)
\end{answer}

\begin{problem}
两部不同的长篇小说各由第一，二，三，四卷组成. 每卷 1 本，共 8 本. 将它们任意地排成一排，左边 4 本恰好都属于同一部小说的概率是\fillinblank{}. （结果用分数表示）
\end{problem}

\begin{answer}
\(\frac{1}{35}\)
\end{answer}

\begin{problem}
如果一条直线与一个平面垂直，那么，称此直线与平面构成一个“正交线面对”. 在一个正方体中，由两个顶点确定的直线与含有四个顶点的平面构成的“正交线面对”的个数是\fillinblank{}.
\end{problem}

\begin{answer}
\(36\)
\end{answer}

\begin{problem}
若曲线 \( y^{2} = |x| + 1 \) 与直线 \( y = kx + b \) 没有公共点，则 \(k\)，\(b\) 分别应满足的条件是\fillinblank{}.
\end{problem}

\begin{answer}
\(k=0\)，\(b\in(-1,1)\)
\end{answer}

\begin{problem}
三个同学对问题“关于 \(x\) 的不等式 \(x^{2} + 25 + |x^{3} - 5x^{2}| \geqslant ax\) 在 \([1, 12]\) 上恒成立，求实数 \(a\) 的取值范围”提出各自的解题思路. 甲说：“只须不等式左边的最小值不小于右边的最大值”. 乙说：“把不等式变形为左边含变量 \(x\) 的函数，右边仅含常数，求函数的最值”. 丙说：“把不等式两边看成关于 \(x\) 的函数，作出函数图象”. 参考上述解题思路，你认为他们所讨论的问题的正确结论，即 \(a\) 的取值范围是\fillinblank{}.
\end{problem}

\begin{answer}
\(a\in(-\infty,10]\)
\end{answer}

\section{单选题}
\begin{problem}
如图，在平行四边形 \(ABCD\) 中，下列结论中错误的是（\quad）

\bitmapfigure[max width=0.40\linewidth]{img/2006/shanghai_science/q13_fig1.png}

\choices
    {\(\overrightarrow{AB} = \overrightarrow{DC}\)}
    {\(\overrightarrow{AD} +\overrightarrow{AB} = \overrightarrow{AC}\)}
    {\(\overrightarrow{AB} -\overrightarrow{AD} = \overrightarrow{BD}\)}
    {\(\overrightarrow{AD} +\overrightarrow{CB} = \overrightarrow{0}\)}
\end{problem}

\begin{answer}
C
\end{answer}

\begin{problem}
若空间中有四个点，则“这四个点中有三点在同一条直线上”是“这四个点在同一个平面上”的（\quad）

\choices
    {充分非必要条件}
    {必要非充分条件}
    {充分必要条件}
    {既非充分又非必要条件}
\end{problem}

\begin{answer}
A
\end{answer}

\begin{problem}
若关于 \(x\) 的不等式 \((1 + k^2)x \leqslant k^4 + 4\) 的解集是 \(M\)，则对任意实数 \(k\)，总有（\quad）

\choices
    {\(2 \in M, 0 \in M\)}
    {\(2 \notin M\)，\(0 \notin M\)}
    {\(2 \in M\)，\(0 \notin M\)}
    {\(2 \notin M\)，\(0 \in M\)}
\end{problem}

\begin{answer}
A
\end{answer}

\begin{problem}
如图，平面中两条直线 \(l_{1}\) 和 \(l_{2}\) 相交于点 \(O\)，对于平面上任意一点 \(M\)，若 \(p\)，\(q\) 分别是 \(M\) 到直线 \(l_{1}\) 和 \(l_{2}\) 的距离，则称有序非负实数对 \((p, q)\) 是点 \(M\) 的“距离坐标”. 已知常数 \(p \geqslant 0\)，\(q \geqslant 0\)，给出下列命题：

\begin{circlelist}
\item 若 \(p = q = 0\)，则“距离坐标”为 \((0, 0)\) 的点有且仅有 1 个；
\item 若 \(pq = 0\)，且 \(p + q \neq 0\)，则“距离坐标”为 \((p, q)\) 的点有且仅有 2 个；
\item 若 \(pq \neq 0\)，则“距离坐标”为 \((p, q)\) 的点有且仅有 4 个.
\end{circlelist}

上述命题中，正确命题的个数是（\quad）

\bitmapinlinefigure[0.45\linewidth][max width=0.40\linewidth]{img/2006/shanghai_science/q16_fig1.png}

\choices
    {0}
    {1}
    {2}
    {3}
\end{problem}

\begin{answer}
D
\end{answer}

\section{解答题}
\begin{problem}
求函数 \( y = 2\cos \left( x + \frac{\pi}{4} \right) \cos \left( x - \frac{\pi}{4} \right) + \sqrt{3} \sin 2x \) 的值域和最小正周期.
\end{problem}

\begin{solution}
值域为 \([-2,2]\)，最小正周期为 \(\pi\).
\end{solution}

\begin{problem}
如图，当甲船位于 \(A\) 处时获悉，在其正东方向相距 20 海里的 \(B\) 处有一艘渔船遇险等待营救. 甲船立即前往救援，同时把消息告知在甲船的南偏西 \(30^{\circ}\)，相距 10 海里 \(C\) 处的乙船，试问乙船应朝北偏东多少度的方向沿直线前往 \(B\) 处救援. （角度精确到 \(1^{\circ}\)）

\bitmapfigure[max width=0.44\linewidth]{img/2006/shanghai_science/q18_fig1.png}
\end{problem}

\begin{solution}
乙船应朝北偏东 \(71^{\circ}\) 的方向沿直线前往 \(B\) 处救援.
\end{solution}

\begin{problem}
在四棱锥 \(P - ABCD\) 中，底面是边长为 2 的菱形. \(\angle DAB = 60^{\circ}\)，对角线 \(AC\) 与 \(BD\) 相交于点 \(O\)，\(PO \perp\) 平面 \(ABCD\)，\(PB\) 与平面 \(ABCD\) 所成的角为 \(60^{\circ}\).

\begin{enumerate}
\item 求四棱锥 \(P-ABCD\) 的体积；
\item 若 \(E\) 是 \(PB\) 的中点，求异面直线 \(DE\) 与 \(PA\) 所成角的大小. （结果用反三角函数值表示）
\end{enumerate}

\bitmapinlinefigure[0.6\linewidth][max width=0.38\linewidth]{img/2006/shanghai_science/q19_fig1.png}
\end{problem}

\begin{solution}
\begin{enumerate}
\item 四棱锥 \(P-ABCD\) 的体积为 \(2\).
\item 异面直线 \(DE\) 与 \(PA\) 所成角的大小为 \(\arccos\frac{\sqrt{2}}{4}\).
\end{enumerate}
\end{solution}

\begin{problem}
在平面直角坐标系 \(xOy\) 中，直线 \(l\) 与抛物线 \(y^{2} = 2x\) 相交于 \(A\)，\(B\) 两点.

\begin{enumerate}
\item 求证：“如果直线 \(l\) 过点 \(T(3,0)\)，那么 \(\overrightarrow{OA} \cdot \overrightarrow{OB} = 3\)” 是真命题；
\item 写出第一问中命题的逆命题，判断它是真命题还是假命题，并说明理由.
\end{enumerate}
\end{problem}

\begin{solution}
\begin{enumerate}
\item 第一问中的命题为真命题.
\item 逆命题为“如果 \(\overrightarrow{OA}\cdot\overrightarrow{OB}=3\)，那么直线 \(l\) 过点 \(T(3,0)\)”，该逆命题为假命题.
\end{enumerate}
\end{solution}

\begin{problem}
已知有穷数列 \(\{a_{n}\}\) 共有 \(2k\) 项（整数 \(k \geqslant 2\)），首项 \(a_{1} = 2\). 设该数列的前 \(n\) 项和为 \(S_{n}\)，且 \(a_{n+1} = (a-1)S_{n} + 2\)（\(n = 1\)，\(2\)，\(\cdots\)，\(2k-1\)）. 其中常数 \(a > 1\).

\begin{enumerate}
\item 求证：数列 \( \{a_{n}\} \) 是等比数列；
\item 若 \(a = 2^{\frac{2}{2k-1}}\)，数列 \(\{b_n\}\) 满足 \(b_n = \frac{1}{n}\log_2(a_1a_2\cdots a_n)\)（\(n = 1\)，\(2\)，\(\cdots\)，\(2k\)），求数列 \(\{b_n\}\) 的通项公式；
\item 若第二问中的数列 \(\{b_n\}\) 满足不等式 \(\left|b_1 - \frac{3}{2}\right| + \left|b_2 - \frac{3}{2}\right| + \cdots + \left|b_{2k-1} - \frac{3}{2}\right| + \left|b_{2k} - \frac{3}{2}\right| \leqslant 4\)，求 \(k\) 的值.
\end{enumerate}
\end{problem}

\begin{solution}
\begin{enumerate}
\item \(\{a_n\}\) 是等比数列.
\item \(b_n=1+\frac{n-1}{2k-1}\).
\item \(k=2,3,4,5,6,7\).
\end{enumerate}
\end{solution}

\begin{problem}
已知函数 \( y = x + \frac{a}{x} \) 有如下性质：如果常数 \( a > 0 \)，那么该函数在 \((0, \sqrt{a}]\) 上是减函数，在 \([\sqrt{a}, +\infty)\) 上是增函数.

\begin{enumerate}
\item 如果函数 \( y = x + \frac{2^{b}}{x} \)（\(x > 0\)）的值域为 \([6, +\infty)\)，求 \( b \) 的值；
\item 研究函数 \( y = x^{2} + \frac{c}{x^{2}} \) （常数 \( c > 0 \)） 在定义域内的单调性，并说明理由；
\item 对函数 \( y = x + \frac{a}{x} \) 和 \( y = x^{2} + \frac{a}{x^{2}} \) （常数 \( a > 0 \)） 作出推广，使它们都是你所推广的函数的情况. 研究推广后的函数的单调性（只须写出结论，不必证明），并求函数 \( F(x) = \left(x^{2} + \frac{1}{x}\right)^{n} + \left(\frac{1}{x^{2}} + x\right)^{n} \) （\( n \) 是正整数）在区间 \( \left[\frac{1}{2}, 2\right] \) 上的最大值和最小值. （可利用你的研究结论）
\end{enumerate}
\end{problem}

\begin{solution}
\begin{enumerate}
\item \(b=\log_2 9\).
\item 在 \(( -\infty,-\sqrt[4]{c})\) 和 \((0,\sqrt[4]{c})\) 上为减函数，在 \(( -\sqrt[4]{c},0)\) 和 \((\sqrt[4]{c},+\infty)\) 上为增函数.
\item 可推广为 \(y=x^n+\frac{a}{x^n}\)（\(a>0\)，\(n\) 为正整数）. 令 \(r=\sqrt[2n]{a}\)：当 \(n\) 为奇数时，在 \((0,r)\) 上为减函数、在 \((r,+\infty)\) 上为增函数、在 \(( -\infty,-r)\) 上为增函数、在 \((-r,0)\) 上为减函数；当 \(n\) 为偶数时，在 \((0,r)\) 上为减函数、在 \((r,+\infty)\) 上为增函数、在 \(( -\infty,-r)\) 上为减函数、在 \((-r,0)\) 上为增函数. \(F(x)\) 的最大值为 \(\left(\frac{9}{2}\right)^n+\left(\frac{9}{4}\right)^n\)，最小值为 \(2^{n+1}\).
\end{enumerate}
\end{solution}
